% ===========================================================================
% INTRODUCTION --- argument skeleton (write body after Method + Evaluation).
% Follows the introduction protocol (skill S13): problem -> why public weights
% do not remove private assets -> gap -> tension -> method -> contributions ->
% evidence summary -> limitation boundary. Contribution bullets must be atomic
% and each independently defensible; do NOT fuse a proven construction with an
% unproven security implication.
% ===========================================================================
\section{Introduction}
\label{sec:intro}

% ---- CENTRAL QUESTION -------------------------------------------------------
% Can a deployment serve a PUBLIC-weight LLM (and adapt it via LoRA) on an
% UNTRUSTED accelerator so that private inputs, intermediate states, KV-cache,
% and LoRA adapters/gradients are not exposed in cleartext to the accelerator,
% while preserving the model's computation at acceptable cost -- and exactly
% which of those protections can we actually substantiate?

% ---- ARGUMENT PLAN (fill prose last) ---------------------------------------
% (1) Deployment problem: public base weights (Llama/Qwen) are downloadable, but
%     the *runtime* assets are private: user prompts, activations, KV-cache,
%     and privately fine-tuned LoRA adapters + their training data/gradients.
% (2) Public weights do NOT eliminate private runtime assets: activation-
%     inversion / embedding-inversion / alignment attacks recover inputs from
%     exposed activations; gradients leak training data. \citeneeded{activation inversion (EIA)}\citeneeded{deep leakage from gradients}
% (3) Gap: cryptographic private inference (FHE/MPC) is exact but expensive on
%     nonlinearities; TEE-only execution wastes the accelerator; prior
%     obfuscation (STIP/ObfuscaTune/Amulet-style) targets different asset sets,
%     TCBs, nonlinear handling, and mostly non-autoregressive or non-training
%     settings. \citeneeded{STIP}\citeneeded{ObfuscaTune}\citeneeded{Amulet}\citeneeded{Slalom}
% (4) Central tension: masking that is cheap and linear-algebra-friendly leaves
%     algebraic INVARIANTS visible (norms, Gram, permutation-invariant stats);
%     hiding those requires either trusted nonlinear work or heavier crypto.
% (5) High-level method: right-action masks + boundary additive pads let the
%     untrusted accelerator do the heavy linear algebra on transformed tensors;
%     a small trusted component holds secrets, does selected nonlinear/boundary
%     ops, recovery, sampling, and auditing (Section~\ref{sec:method}).
%
% ---- CONTRIBUTIONS (atomic; each maps to ledger IDs -- keep in sync) --------
% C1. A masked-execution construction for public-weight decoder LLMs whose
%     linear, attention (incl. RoPE/GQA/MQA), and KV-cache steps are
%     ALGEBRAICALLY correctness-preserving under stated invertibility/tying
%     conditions.                          -> ledger C-COR-01..06
% C2. A LoRA forward/backward masking scheme that preserves the adapter update
%     for SGD-type optimizers, with an explicit optimizer-compatibility
%     boundary.                            -> ledger C-LORA-*, C-OPT-*
% C3. A leakage characterization that states precisely which invariants remain
%     visible (permutation-invariant statistics, norms, padded rank, shapes),
%     plus an EMPIRICAL attack evaluation under a stated attacker model -- NOT a
%     cryptographic security proof.        -> ledger C-SEC-*
% C4. A cost characterization separating measured wall-clock (CPU emulation) and
%     trusted-boundary-call accounting / operation-count projections; and a
%     real-model functional validation (Qwen2.5-0.5B, CPU, float32, greedy).
%                                          -> ledger C-EFF-*, C-REAL-*
%
% Each bullet above is deliberately split so that "correctness" claims never
% silently import a "security" claim (skill S13).
%
% ---- EVIDENCE SUMMARY (1 paragraph, calibrated) -----------------------------
% \TODO{Summarize: algebraic proofs; float64 verification magnitudes; real-model
%   greedy token-match=1.0 on Qwen-0.5B CPU; attack proxy outcomes; costs are
%   CPU-emulation / projections, NOT GPU/TEE wall-clock.}
%
% ---- LIMITATION BOUNDARY (state up front) -----------------------------------
% \TODO{State plainly: no cryptographic/semantic-indistinguishability claim; no
%   real TEE/GPU deployment measured; correctness validated for GREEDY decoding
%   only; RMSNorm currently uses trusted recompute in the wired path; padded
%   rank is visible; norms/permutation-invariants leak.}
%
% ---- REVIEWER RISKS ---------------------------------------------------------
% \risk{Overclaiming security from algebraic correctness.}
% \risk{Novelty vs STIP/ObfuscaTune/Amulet must be crisp on asset set + TCB.}
% \risk{"Public-weight" framing must not imply nothing is secret.}
\TODO{Write Introduction prose after Method and Evaluation stabilize; import calibrated numbers from claim\_ledger.md only.}
